\documentclass[11pt]{article}


\usepackage[margin=1in]{geometry}
\usepackage{booktabs}
\usepackage{multirow}
\usepackage{caption}

\captionsetup{
  font=small,
  labelfont=bf
}

\begin{document}

% This document is meant to be overwritten from the notebook with
% write_model_performance_latex_document(paper_table, ...).
% The table formatting matches the paper; the values should come from
% your generated model results.

\begin{table}[ht]
\centering
\caption{CP and FV Model performance compared to baseline.}
\label{tab:cp-fv-performance}
\begin{tabular}{llrrrr}
\toprule
 & Metric & FV Baseline & FV Model & CP Baseline & CP Model \\
\midrule
\multirow{4}{*}{Random} & Accuracy & 0.708 & 0.769 & 0.932 & 0.946 \\
 & AUC & 0.622 & 0.791 & 0.726 & 0.906 \\
 & PPV & 0.709 & 0.769 & 0.935 & 0.947 \\
 & NPV & 0.512 & 0.765 & 0.454 & 0.905 \\
\midrule
\multirow{4}{*}{Temporal} & Accuracy & 0.694 & 0.760 & 0.934 & 0.948 \\
 & AUC & 0.629 & 0.787 & 0.713 & 0.900 \\
 & PPV & 0.696 & 0.759 & 0.939 & 0.949 \\
 & NPV & 0.446 & 0.762 & 0.423 & 0.903 \\
\midrule
\multirow{4}{*}{Player} & Accuracy & 0.730 & 0.789 & 0.936 & 0.950 \\
 & AUC & 0.607 & 0.789 & 0.755 & 0.913 \\
 & PPV & 0.730 & 0.792 & 0.941 & 0.951 \\
 & NPV & 0.482 & 0.773 & 0.461 & 0.933 \\
\bottomrule
\end{tabular}
\end{table}

% To overwrite this document with your own model output:
%
% 1. In your notebook, run:
%
%    from ufa import write_model_performance_latex_document
%
%    write_model_performance_latex_document(
%        paper_table,
%        "../docs/model_performance_table.tex",
%    )
%
% 2. Rebuild the PDF in VS Code with the pdflatex direct recipe.

\end{document}
